% The 20-Round Barrier for Table-Based Preimage Attacks on SHA-256
% (companion note to the 19-round preimage paper)
%
% Compile: pdflatex paper_r20_barrier.tex  (twice for cross-references)

\documentclass[11pt,a4paper]{article}

\usepackage{amsmath,amssymb,amsthm}
\usepackage{booktabs}
\usepackage{xcolor}
\usepackage[hidelinks]{hyperref}
\usepackage{geometry}
\geometry{margin=1in}
\usepackage[protrusion=true,expansion=false]{microtype}

\newtheorem{theorem}{Theorem}
\newtheorem{corollary}{Corollary}
\newtheorem{definition}{Definition}
\newtheorem{lemma}{Lemma}
\newtheorem{proposition}{Proposition}
\newtheorem{remark}{Remark}

\newcommand{\Sig}[1]{\Sigma_{#1}}
\newcommand{\sig}[1]{\sigma_{#1}}
\newcommand{\Maj}{\mathrm{Maj}}
\newcommand{\Ch}{\mathrm{Ch}}
\newcommand{\ROTR}[2]{\mathrm{ROTR}^{#1}(#2)}
\newcommand{\SHR}[2]{\mathrm{SHR}^{#1}(#2)}
\newcommand{\mmod}{\bmod\,2^{32}}


\begin{document}

\title{\textbf{The 20-Round Barrier for Table-Based\\
Preimage Attacks on SHA-256}}

\author{Kameldip Singh Basra \\[4pt]
        \small Independent Researcher \\
        \small \texttt{kameldipbasra@gmail.com} \\
        \small ORCID: \href{https://orcid.org/0009-0001-0977-4614}{0009-0001-0977-4614}}
\date{September 2026}

\maketitle

\begin{abstract}
This is a companion note to~\cite{Basra2026R19}, which gives a practical
preimage attack on the 19-round SHA-256 compression function built on a
precomputed table inverting $u\mapsto\sigma_0(u)-u$.  Here we ask why the same construction
does not reach 20 rounds, and characterise the barrier quantitatively rather
than merely reporting it.

The $R{=}19$ attack is shown not to schedule its constraints acyclically: it
breaks a cycle by freezing $W_9$ provisionally, iterating, and paying one
self-consistency check.  Promoting $a_{11}$ to an unknown absorbed by the
fourth constraint gives $R{=}20$ the same one-word surplus, and the true
$R{=}20$ solution is verified to be a fixed point of all four absorber maps---so
the architecture transplants in form.  It does not transplant in efficacy:
measured with a single harness, the $R{=}19$ two-map iteration converges at
$5.34\times10^{-5}$ per $a_0$ ($1{,}790$ of $3.36\times10^{7}$) while the
$R{=}20$ four-map iteration yields $0$ of $2.00\times10^{9}$
($\le1.51\times10^{-9}$ at 95\%), below the $2^{-29}$ break-even.  The cause is
dimensional: the iterated state grows from $32$ to $64$ bits.

We then give the structural reason.  Each free state word $a_k$ enters
exactly two recovered message words linearly, $W_k$ with coefficient $+1$ and
$W_{k+8}$ with coefficient $-1$, so constraint C$_j$ absorbs $a_{j+1}$ through
the one global $\sigma_0(u){-}u$ table for \emph{every} $j$, C3$\to a_4$
included.  At $R{=}19$ dependence among the absorbed words is
full-avalanche only in the forward direction and every feedback edge
passes through Maj/Ch alone (at most $1.9$ bits of a constraint change per
flipped bit); at $R{=}20$, in the frame that absorbs $a_4$, exactly one
feedback edge is heavy, $a_4\to$C0 through $W_9$ ($12.2$ bits), and in the
frame that absorbs $a_{11}$ three are; either way the heavy dependency
graph acquires its first cycle.  The $R{=}19$ iteration works because it
cuts its cycle at a mild edge; a heavy cycle has no mild cut, and any
$32$-bit cut of it behaves, as far as we can measure, like a random map.  The
obstruction is $\Sigma_0(a_4)$ subtracted directly inside $W_9$, which no
context choice removes; it is the reach of the round function, not of the
frame: the window is all-mild exactly when $R\le19$, and $R{=}21$ has three
heavy feedback edges.  We give the corrected inventory of which absorptions
admit a global table, show that with the tables we know the assignment is
forced up to a choice between two frames, both of which close the heavy
cycle, and record the resulting cost: $2^{64}$ per context-equivalent,
i.e.\ $2^{32}$ times the $R{=}19$ per-context cost, for this attack
family.  We record a second global
table (the $a_0$ absorber, $\Gamma(a_0)=a_0+g(a_0)$, coverage $63.214\%$
against $1-e^{-1}=63.212\%$), and prove that
$\operatorname{rank}(\sig{0}\oplus I)=31$ gives a unique mask
$\lambda=\texttt{0xa8a42fe4}$ with $\lambda\cdot(\sig{0}(W)\oplus W)=0$ for
every $W$ (verified exhaustively over all $2^{32}$ inputs), which modular
subtraction degrades to a bias of only $2^{-15.8}$.

Because these are properties of the parameterisations we tried rather than
proofs of impossibility, we also build a reduced-width scale model in which the
whole inner space is enumerable, and report exhaustive measurements there.
\textbf{None of this is an impossibility result}, and we are explicit
throughout about which claims are algebraic and which are statistical models
that remain untested at full width.  Full SHA-256 (64 rounds) is not threatened
by any of these results.

\medskip
\noindent\textbf{Keywords:} SHA-256, preimage attack, reduced-round,
round-20 barrier, absorber tables, reduced-width scale model.
\end{abstract}

\section{Scope and Relation to the \texorpdfstring{$R=19$}{R=19} Paper}

We assume the attack of~\cite{Basra2026R19} throughout: a backward chain from
the target fixes $a_{12},\ldots,a_{19}$ (for $R{=}20$), the context words
$a_4,\ldots,a_{10}$ are freely chosen (one of the two admissible frames,
\S\ref{sec:frames}), and the remaining unknowns are resolved
against the schedule-consistency constraints C0--C3 at $W_{16},\ldots,W_{19}$
using a $16$\,GB global table inverting $u\mapsto\sig{0}(u)-u$.

\paragraph{Notation.}
We follow~\cite{Basra2026R19}: $T_1^{(r)}$, $T_2^{(r)}$ and $e_r$ are the
round quantities of round $r$; $g(a_0) = W_1 - a_1$ is the part of $W_1$
that depends only on $a_0$, as defined there; $D(a_3)$ is the
$a_3$-dependent part of $W_{10}$ and $F_{12}$, $F_{23}$ the
$(a_0,a_1)$- and $(a_0,a_1,a_2)$-dependent parts of $W_2$ and $W_3$, as in
its Proposition~3; and $c_0 = W_0 - a_0$ is the round-0 constant (its
$C_0$).  Constants that depend only on the context and the target are
written $\mathrm{CONST}$ with a subscript.  $R_r$ for $r=16,\ldots,19$
denotes the residual of constraint C$_{r-16}$: the message word $W_r$
recovered from the state minus its schedule value, zero exactly when the
constraint holds.  $H(a_{11})$ denotes the
right-hand side of C3 as a function of $a_{11}$ alone and $h_1(a_{11})$ the
$a_{11}$-dependent part of the C1 target.  A \emph{frame} is a choice of
which state words are context and which are absorbed
(\S\ref{sec:frames}).

\paragraph{Scope.}
``Barrier'' here is specific to this attack family.  Meet-in-the-middle
and biclique preimage attacks reach 43 and 45 steps at costs near
$2^{256}$~\cite{AokiGuo2009,KRS2012}, and Zaikin~\cite{Zaikin2026} has
inverted the 19-round compression function for a fixed digest by
parameterized Cube-and-Conquer; nothing here bears on those results.

Nothing here modifies the $R{=}19$ result.  The material is separated because
it is exploratory: it consists largely of negative results, and its value is in
narrowing where a $R{=}20$ attack could come from, not in delivering one.

\section{The \texorpdfstring{$R=20$}{R=20} Constraint System}
\label{sec:r20}
%-----------------------------------------------------------------------

The backward chain at $R = 20$ recovers $a_{12},\ldots,a_{19}$ but
not $a_{11}$.  The schedule grows to four constraints
$W_{16}$--$W_{19}$, with $a_{11}$ as an additional unknown.

\textbf{Reference-conditioned result.}  When the context including $a_{11}$
is taken from a known preimage of the target, the $W_{17}$ identity
gives $a_2$ as an explicit function of $a_3$ (Lemma analogous to
\cite[Prop.~3]{Basra2026R19}), and an $h$-table $h[W_{12}(a_{11})]
\to a_{11}$ enables an $O(2^{32})$ $a_3$ sweep.  This recovered a
verified R=20 preimage in 227\,s of sweep time plus $\approx 22$
minutes for the $h$-table build.  However, the context was derived
from a known preimage: this is a second preimage computed with knowledge
of the first, not a preimage attack in the standard sense.

\textbf{Why the R=20 preimage attack is not established.}  The $h$-table
inversion requires $W_{12}^{\mathrm{req}} = W_{19}^{\mathrm{bc}} -
\sig{1}(W_{17}^{\mathrm{bc}}) - \sig{0}(W_4) - W_3$, where
$W_{17}^{\mathrm{bc}}$ carries $a_{11}$ through $e_{13}$
(which contains $\Maj(a_{12},a_{11},a_{10})$).  Thus $a_{11}$
is needed to compute the quantity used to look up $a_{11}$: a circular
dependency.  No algebraic link that breaks this circularity has been
found.  Exhaustive search over all binary linear combinations of the
four residual equations $R_{16}$--$R_{19}$ found no combination that
isolates any of the three unknowns $(a_2, a_3, a_{11})$.

\textbf{Quantitative complexity gap.}
An equivalent view using the $\sigma_0$-chain directly confirms the $2^{32}$ difficulty increase.
In the classical accounting, refined in \S\ref{sec:frames}, treat $a_4,\ldots,a_{11}$ as freely chosen context (the backward chain fixes $a_{12},\ldots,a_{19}$).
The chain C0--C2 then proceeds identically to the $R=19$ case, but the schedule now includes $W_{19}$,
introducing a fourth constraint~(C3):
\[
  \sig{0}\!\bigl(W_4^{\mathrm{round}}\bigr)
  = W_{19}^{\mathrm{bc}} - \sig{1}(W_{17}^{\mathrm{bc}}) - W_{12}^{\mathrm{base}} - W_3,
\]
where $W_4^{\mathrm{round}}$ is fully determined by $a_0,\ldots,a_3$ once C0--C2 converge and
$W_{12}^{\mathrm{base}}$ is determined by context.
Because $\sig{0}$ is a bijection, $W_4^{\mathrm{round}}$ must equal the unique $\sig{0}^{-1}(C3_{\mathrm{tgt}})$.
C3 therefore contributes an additional 32-bit equality.  Under the present
construction we model it as approximately uniform and independent of C0--C2, so
that it passes with probability $\approx 2^{-32}$; we stress that this is a
statistical model which has not been established directly at full width.

Exhaustive experiments in a reduced-width scale model (\S\ref{sec:redwidth})
bear on the model but do not confirm it at $w=32$.  Measuring
$P(\text{C3}\mid \text{C0},\text{C1},\text{C2})$ against the independent
prediction $2^{-w}$ gives a ratio of $1.075\pm0.018$ at $w=4$ and
$1.005\pm0.035$ at $w=5$: consistent with exact independence at the larger
width, with a small excess at the smaller one that may be a narrow-width
artefact.  Any dependence is thus small, but two widths do not establish a
trend, and the behaviour at $w=32$ remains untested.

\textbf{Schedule identity.}
For messages satisfying the $\sigma_0$-chain $W_k = \sig{0}(W_{k+1})$, $k=0,\ldots,3$,
the substitution $\sig{0}(W_{r-15}) = W_{r-16}$ reduces the schedule recurrence for $r\in\{16,\ldots,19\}$ to
\[
  W_r = \sig{1}(W_{r-2}) + W_{r-7} + 2\,W_{r-16}.
\]
From this identity and $W_3 = \sig{0}(W_4)$ one computes directly
$C3_{\mathrm{tgt}} = W_{19} - \sig{1}(W_{17}) - W_{12} - W_3 = \sig{0}(W_4) = W_3$,
confirming that the C3 check reduces to $W_4^{\mathrm{round}} = W_4^{\mathrm{chain}}$ (verified numerically).
The same identity provides a concise derivation of the Nine-Step Lemma: at $r=16$,
$\Delta W_{16} = \Delta W_{9} + \Delta W_0 = -\delta + \delta = 0$, and the doubled term
$2W_{r-16}$ ensures that \emph{any} perturbation $\Delta W_1 = \varepsilon$ would
introduce $\sig{0}(\varepsilon)$ into $\Delta W_{16}$ (breaking cancellation), which is
why the pair $(0,9)$ is the unique minimal solution.

\textbf{Orbit structure of $\boldsymbol{\sigma_0}$.}
Viewed as a $\mathbb{F}_2$-linear bijection on $(\mathbb{Z}/2\mathbb{Z})^{32}$, the map $\sigma_0$
has order $49{,}146 = 2\cdot 3\cdot 8191$ where $8191 = 2^{13}-1$ is a Mersenne prime.
Its only short orbits are the two fixed points $\{0,\,\texttt{0x27f42515}\}$
and the unique $2$-cycle $\{\texttt{0x0b3b7ef9} \leftrightarrow \texttt{0x2ccf5bec}\}$;
in particular, no period-3 or period-4 elements exist.
Consequently any $\sigma_0$-chain of length four ($W_0,\ldots,W_4$) lies in an orbit of
length $49{,}146$ unless $W_4$ is one of the four exceptional values above, so
no ``chain-wrap'' shortcut of the form $W_0 = W_4$ arises generically.

Under the variable ordering and single-representative lookup construction used
here, C3 remains an additional 32-bit consistency condition applied after
C0--C2 have been solved.  We emphasise what is and is not established.  The
algebra shows C3 imposes a further 32-bit equality; it does not show that every
algebraic arrangement must pay $2^{32}$ for it.  No absorption of C3 into the
construction, and no lower-width decomposition of it, is currently known---but
the negative results below are properties of the parameterisations we tried,
not proofs of impossibility.  Indeed one apparent obstruction at R=20, the
$a_1\!\leftrightarrow\!a_2$ cycle, dissolved entirely under a change of
variable (\S\ref{sec:ureparam}), which is direct evidence that such
obstructions can be artefacts of representation.  Treating C3 as an independent
filter gives the cost in one convention, which we use throughout.  An
$R{=}19$ attempt costs $2^{32}$ table lookups and succeeds with probability
$\mu_{R=19}$, measured at $1.25\times10^{-2}$ and projected at
$6\times10^{-3}$--$1.3\times10^{-2}$~\cite[\S5.6]{Basra2026R19}.  An
$R{=}20$ attempt in the classic harness guesses one further word and costs
$2^{64}$, so we say the barrier is \emph{$2^{64}$ per context-equivalent}.
The expected number of $R{=}19$-sized contexts is then
$\mu_{R=19}^{-1}\times 2^{32}\approx 2^{39.4}$ at the projected rate, and
the expected total work $\mu_{R=19}^{-1}\times 2^{64}\approx 2^{71.4}$.
The factor of $2^{32}$ between the two round counts is what matters and is
independent of $\mu_{R=19}$; the absolute figures inherit its uncertainty.

\textbf{GPU diagnostic.}
A 165-context GPU run of the $R=20$ $\sigma_0$-chain on an H100 (batch~$=2^{26}$,
K\textsubscript{seeds}$=4$; script \texttt{r20\_oracle\_free.py}, 19 June
2026, a separate run from the 165-context $R{=}19$ run reported
in~\cite{Basra2026R19}) recorded 4{,}166 $W_9$-lo-pass events and zero
hi-pass events.  Neither this log nor that of the CPU sweep below is
included in the artefact package.
Under the null hypothesis of uniform residuals, the expected hi count is $\mu =
4166/65536 \approx 0.064$; the probability of zero hi events is $e^{-\mu}\approx 0.94$,
so the absence of hi events is unsurprising.  With an expectation
below $0.1$, however, observing zero carries almost no information: the result
is \emph{consistent with} a uniform $2^{-32}$ filter but does not test it.
We attribute the difficulty increase over $R=19$ to the additional
C3 equality rather than to any structural bias in $W_9$; the attribution of the
full $2^{32}$ factor to C3 rests on the uniformity model stated above, which
this run is too small to test.
A supplementary CPU sweep (10 contexts, batch~$=2^{24}$, two workers;
script \texttt{r20\_corr\_parallel.py}, 20 June 2026) recorded 254
$W_9$-lo-pass events and zero C3-pass events.  With zero events the
one-sided 95\% Poisson bound is $3.0$, so $P(\mathrm{C3}\mid\mathrm{lo})
\le 3/254$ and the enhancement over a uniform $2^{-32}$ filter is at most
$3\cdot 2^{32}/254 = 2^{25.6}$.  This too is consistent with a uniform
filter and far too small a sample to test independence or estimate its
probability.  The structural claim---that C3 imposes an additional 32-bit
equality---rests on the algebra, not on these counts; the probabilistic claim
should be read as untested.

\subsection{A Reduced-Width Scale Model}
\label{sec:redwidth}

Claims about C3's independence, and about whether the chained solver is lossy,
are untestable at full width: the events are too rare to sample and the
solution set cannot be enumerated.  We therefore build a scale model of the
round function and message schedule at word width $w$, where the inner space
$(a_0,a_1,a_2,a_3,a_{11})$ is small enough to enumerate exhaustively.
Rotation and shift amounts are scaled from SHA-256's by $w/32$ and rounded; $K$
and the IV are truncated.  The construction preserves the shape of the mixing
and is not proposed as a standard.

Two fidelity checks are passed.  At $w=32$ the scaled constants reduce to
SHA-256's exactly, and the model reproduces SHA-256's state evolution
word for word.  The solution count per context matches the degrees-of-freedom
prediction: five unknown words less four constraints leaves one surplus word,
hence $2^w$ solutions, against $19.2$ measured at $w=4$ ($2^4=16$) and $32.0$
at $w=5$ ($2^5=32$).

The independence measurement is reported in \S\ref{sec:r20}.  The model also
shows that the chained fixed-point architecture is intrinsically lossy rather
than merely slow.  A solver using one representative per table entry recovers
$1.3\%$, $0\%$ and $0\%$ of the true solutions at $w=4,5,6$, while one
branching over all roots recovers $23.4\%$, $6.8\%$ and $0.7\%$.  Raising the
beam from 64 to 1024 and the iteration count from 6 to 20 changed coverage not
at all, so the loss is structural rather than a truncation effect.  Tables
should therefore be relations, not functions.  This monotone collapse, resting
on 36, 13 and 2 recovered solutions, is the model's robust finding.

We deliberately do \emph{not} extrapolate these rates to $w=32$.  The
per-$a_0$ success rate falls by $4.15\times$ from $w=4$ to $w=5$ and by
$8.67\times$ from $w=5$ to $w=6$, so the point estimates suggest an
accelerating rather than log-linear decay; but the $w=6$ figure rests on two
events, whose exact Poisson interval puts the second ratio anywhere in
$[2.40\times, 71.6\times]$, a range that contains the first.  The data
therefore constrain neither the functional form nor the rate, while a $5\%$
change in the fitted rate, together with the accompanying shift of the
$w=5$ base point, moves a 27-bit extrapolation by a factor of $4.5$
($3.8$ from the rate alone).  The scale model's value lies in its
exhaustive in-range comparisons, not in extrapolated rates.

\subsection{Absorber Reparameterisation and the \texorpdfstring{$a_1\!\leftrightarrow\!a_2$}{a1-a2} Coupling}
\label{sec:ureparam}

The four residual equations can each be written in the $\sig{0}(u)-u$ form used
by the R=19 table.  Substituting the round-function expression
$W_{10} = \mathrm{CONST}_{10} - D(a_3) - a_2$ into the $W_{17}$ identity and then
eliminating $a_2$ via $a_2 = W_2 - F_{12}(a_1)$ yields:

\begin{lemma}[Reparameterisation of C1]
\label{lem:ureparam}
Write $A(a_1) := a_1 + g(a_0) + F_{12}(a_1)$, a function of $a_1$ alone once
$a_0$ is fixed.  Then
\[
  \sig{0}(W_2) - W_2 \;=\; h_1(a_{11}) - \mathrm{CONST}_{10} + D(a_3) - A(a_1)
  \;=:\; u ,
\]
so the target $u$ decomposes additively into independent contributions from
$a_{11}$, $a_3$ and $a_1$.
\end{lemma}

Verified exactly on $2000/2000$ random triples $(a_1,a_3,a_{11})$.  Taking $u$ as
the free variable inverts the map: $W_2$ is one table lookup,
$a_1 = A^{-1}(\kappa-u)$ with $\kappa := h_1(a_{11}) - \mathrm{CONST}_{10} +
D(a_3)$ a second, and $a_2 = W_2 - F_{12}(a_1)$.  Hence for fixed $(a_3,a_{11})$ each $u$
determines \emph{both} $a_1$ and $a_2$ in $O(1)$.  The apparent circular
dependency between C0 (which needs $a_2$) and C1 (which needs $a_1$) is therefore
an artefact of the parameterisation, not a genuine fixed-point problem; this
explains why naive iteration of that pair never converges.  ($A$ behaves
like a random map: $3$ collisions in $2\times10^{5}$ samples against $4.7$
expected for a random map and $0$ for a bijection, so $A^{-1}$ is a
relation with about $63\%$ coverage and requires a bucketed table.)

Two further reductions place the remaining constraints in table form.  In C2 the
unknown $a_3$ cancels from its own target, since $e_7 = a_3 + T_1^{(7)}$ and
$W_3 = a_3 + F_{23}$:
\[
  \sig{0}(W_3) - W_3 = \bigl(W_{18} - \sig{1}(W_{16}) - W_{11}\bigr) - W_2 - F_{23} + T_1^{(7)} .
\]
In C3, the quantities $e_{16},e_{17},e_{18}$ do not depend on $a_{11}$ (they are
fixed by the backward chain), so $\sig{1}(e_{18})$ and $\Ch(e_{18},e_{17},e_{16})$
are per-context constants and
\[
  W_{19} \;=\; \mathrm{CONST}_{19} - \bigl(a_{11} + T_1^{(15)}\bigr)
\]
is \emph{linear} in $a_{11}$; only $W_{17}$ and $W_{12}$ carry nonlinear
$a_{11}$-dependence into the C3 key.

\subsection{A Context Family that Removes \texorpdfstring{$a_2$}{a2} from C0}
\label{sec:a2free}

The variable $a_2$ reaches the C0 target by exactly two routes: the term
$\Ch(e_8,e_7,e_6)$ with $e_6 = a_2 + T_1^{(6)}$, and $\Maj(a_4,a_3,a_2)$
inside $T_1^{(5)}$.  Both are maskable by context choice.  Setting
$e_8 = \texttt{0xFFFFFFFF}$ (reachable by choosing $a_8$, since
$e_8 = a_4 + a_8 - \Sig{0}(a_7) - \Maj(a_7,a_6,a_5)$) collapses
$\Ch(e_8,e_7,e_6)$ to $e_7$; setting $a_4 = a_3$ collapses
$\Maj(a_4,a_3,a_2)$ to $a_3$.  The second condition ties a context word to
an absorbed unknown, so the context constants would have to be rebuilt for
every candidate $a_3$; it is recorded as a structural observation, not as
an attack step.  Counting distinct C0 targets over several hundred random
$a_2$ per trial: with a random context, or with either condition alone,
essentially every $a_2$ gives a distinct target; with both conditions the
count is exactly one in each of four trials.

Neither condition suffices alone; together they eliminate $a_2$ from C0 entirely,
turning it into a two-input absorber $(a_3,a_{11})\mapsto a_1$.  This is the first
demonstration that the R=20 constraint \emph{structure} is shapeable by context
choice rather than fixed.  It does not by itself reduce the constraint count.

\subsection{The R=19 Iteration Architecture Does Not Transfer Directly}
\label{sec:noniter}

The R=19 attack does not schedule its constraints acyclically: it too has a cycle
(C0 requires $W_9$, which requires $a_2,a_3$, which come from C1/C2, which require
$a_1$ from C0).  It breaks the cycle by freezing $W_9$ at a provisional value,
iterating C1$\leftrightarrow$C2 to a fixed point, and paying one $32$-bit
self-consistency check.  Convergence of that iteration is what satisfies C1 and C2.

Promoting $a_{11}$ from guessed context to an unknown absorbed by C3 gives R=20 the
same one-word surplus ($5$ unknowns $a_0..a_3,a_{11}$ against $4$ constraints,
versus R=19's $4$ against $3$), and we verified that the true R=20 solution is a
fixed point of all four absorber maps.  The architecture therefore transplants in
form.  It does not transplant in efficacy.  Measured with one harness, cold start,
$8$ iterations, no seeding, on the same hardware:

\begin{center}
\begin{tabular}{lrl}
\hline
iteration & convergences / trials & rate per $a_0$ \\
\hline
R=19 control (two-map, C1$\leftrightarrow$C2) & $1{,}790$ / $33{,}554{,}432$ & $5.34\times10^{-5}$ \\
R=20 (four-map, C0--C3) & $0$ / $1{,}996{,}488{,}704$ & $\le1.51\times10^{-9}$ (95\%) \\
\hline
\end{tabular}
\end{center}

The R=19 control gives $5.34\times10^{-5}$ per $a_0$.  The main paper
documents $1.3$--$1.9\times10^{-5}$ per C0 survivor cold-start, i.e.\
$0.8$--$1.2\times10^{-5}$ per $a_0$, and about $3.3\times$ more with
seeding~\cite[\S5.4]{Basra2026R19}; the control therefore sits within a
factor of five of the cold-start figure and close to the seeded one.  The
harnesses differ in iteration cap and deduplication, and the $R{=}20$
conclusion below does not depend on which figure is taken.  The cost
accounting is as follows.  Per context the iterated variant spends
$2^{32}\cdot 8 = 2^{35}$ steps and yields $2^{32}p$ convergences, each
passing the final 32-bit check with probability $2^{-32}$ under the uniform
model, hence $p$ preimages per context and $2^{35}/p$ steps per preimage;
the classic harness spends $2^{32}$ per guessed $a_{11}$ over $2^{32}$
guesses, i.e.\ $2^{64}$ per context-equivalent.  The break-even
convergence rate is therefore $p = 2^{-29}\approx1.86\times10^{-9}$, and
the measured bound lies below it.  We note the margin is modest---the $95\%$
bound is a factor $1.24$ under break-even---so the conclusion is that the
iterated variant does not beat the classic sweep, not that it loses by a wide
factor; a longer run would tighten the bound proportionally.  The structural
estimate is far more pessimistic than the bound: R=19's $5.34\times10^{-5}$ sits
about $2^{15}$ above the $2^{-29}$ random-map baseline for a $32$-bit state, and
carrying that same enhancement to a $64$-bit state (baseline $\approx2^{-61}$)
predicts $p\approx2^{-46}$, some five orders of magnitude below break-even.  The cause is dimensional: at R=19 the absorbed
$a_1$ is computed \emph{outside} the loop, leaving a $32$-bit iterated state
($a_3$, with $a_2$ derived), whereas at R=20 C0 depends on $a_{11}$ and the state
becomes $64$-bit $(a_3,a_{11})$.  Raising the iteration budget does not help, since
convergence probability and cost scale alike.  Section~\ref{sec:heavy} gives the
structural reason the state cannot be held at $32$ bits in \emph{any} frame:
at $R{=}20$ the full-avalanche dependency graph among the absorbed words
acquires a cycle, and a cycle with no Maj/Ch-only edge has no cheap cut.

\subsection{A Second Global Table: the \texorpdfstring{$a_0$}{a0} Absorber}
\label{sec:a0absorber}

The attack sweeps $a_0$ and uses C0 to absorb $a_1$ through the global
$\sig{0}(u)-u$ table.  The two roles can be exchanged.  Using
\[
  W_1 = a_1 + g(a_0), \qquad W_0 = a_0 + c_0, \qquad
  W_9(a_1) = W_9^{(0)} - a_1,
\]
where $W_9^{(0)}$ is $W_9$ evaluated at $a_1=0$, the constraint
$W_{16} = \sig{1}(W_{14}) + W_9 + \sig{0}(W_1) + W_0$ rearranges to give:

\begin{proposition}
\label{prop:a0absorber}
Let $\Gamma(a_0) := a_0 + g(a_0)$ and
$R := W_{16} - \sig{1}(W_{14}) - W_9^{(0)} - c_0$.  Then C0 is
equivalent to
\[
  \Gamma(a_0) \;=\; R + W_1 - \sig{0}(W_1).
\]
\end{proposition}

The point is that $g$ depends only on $a_0$ and the SHA-256 IV---no context
word and no target enters it---so $\Gamma$ is a fixed function and $\Gamma^{-1}$
is a \emph{global} table, built once and valid for every target and context,
exactly like $\sig{0}(u)-u$.  Measured over all $2^{32}$ inputs, $\Gamma$ covers
$63.214\%$ of the range against $1-e^{-1}=63.212\%$ for a random map, so the
table behaves like the one already in use.  Proposition~\ref{prop:a0absorber} is
verified on the three preimages of~\cite[Table~2]{Basra2026R19} by
\texttt{code/verify\_a0\_absorber.py}, which also confirms it is not a tautology:
perturbing $W_{16}$ by $1,\ldots,64$ breaks it in every case.

This does not reduce the attack's cost.  C0 can absorb $a_0$ \emph{or}
$a_1$, but not both: exchanging the roles gives a mirror frame that sweeps
$W_1$ and absorbs $a_0$, with the same number of absorbed words and the
same dependency structure.  What the second table changes is the inventory
of \S\ref{sec:linear}, not the count; why the count cannot be improved is
the edge-weight argument of \S\ref{sec:heavy}.

\subsection{A Linear Identity for \texorpdfstring{$\sigma_0$}{sigma0}, and its Destruction}
\label{sec:leftnull}

Since $\sig{0}$ is $\mathbb{F}_2$-linear, put $L = \sig{0}\oplus I$.  The kernel of
$L$ is the fixed-point set $\{0,\texttt{0x27f42515}\}$, so $\operatorname{rank}L = 31$
and the image is a $31$-dimensional subspace.  Consequently there is a
\emph{unique} nonzero mask annihilating it.

\begin{proposition}
\label{prop:leftnull}
Let $\lambda = \texttt{0xa8a42fe4}$.  Then for every $W\in(\mathbb{Z}/2\mathbb{Z})^{32}$,
\[
  \lambda \cdot \bigl(\sig{0}(W)\oplus W\bigr) = 0 .
\]
\end{proposition}

Verified exhaustively over all $2^{32}$ inputs (zero violations).  This is the left
companion of the known right-kernel constant \texttt{0x27f42515}.

The identity is however destroyed by modular subtraction, which is what the attack
actually uses.  Exhaustively over all $2^{32}$ inputs,
\[
  \Pr_W\bigl[\lambda\cdot(\sig{0}(W) - W) = 1\bigr] = 0.500017190 ,
\]
a bias of $1.72\times10^{-5} = 2^{-15.8}$, against an exact $0$ for the XOR form.
A relation that is identically zero on the $\mathbb{F}_2$-linear part survives
modular subtraction only as a $2^{-15.8}$ bias, which affords no usable filtering.
Together with the table coverage---the $\sig{0}$ table covers $63.37\%$ of
its range, the $\Gamma$ table $63.214\%$ and the C3 $h$-table $63.21\%$,
against $1-e^{-1}=63.212\%$ for a random map---the map
$W\mapsto\sig{0}(W)-W$ is indistinguishable from a random map under every
structural test we have applied.  We conjecture that the interaction
between $\mathbb{F}_2$-linear $\sig{0}$ and modular arithmetic is not
merely an obstacle to proof technique but the mechanism by which R=20
resists.

The R=20 preimage attack remains an open problem.

%-----------------------------------------------------------------------


\section{The Absorber Assignment, and the Heavy-Cycle Barrier}
\label{sec:forced}

This section gives the sharpest statement of the barrier we have, in two
steps: an inventory of which absorptions admit a \emph{global} table, and
a measurement of the dependency structure among the absorbed words that
explains, with one number per edge, why the $R{=}19$ iteration
architecture cannot be carried to $R{=}20$ in either admissible frame.
Neither step is an argument about cycles alone: the $R{=}19$ assignment is
also cyclic and succeeds.  What matters is the \emph{weight} of the edges.

\subsection{Linear Entries and the Global Absorbers}
\label{sec:linear}

Write $T_2(x,y,z)=\Sig{0}(x)+\Maj(x,y,z)$.  In the $a$-only representation
of~\cite{Basra2026R19}, $e_r = a_{r-4} + a_r - T_2(a_{r-1},a_{r-2},a_{r-3})$
and the recovered message word is
\[
  W_r \;=\; a_r - T_2(a_{r-1},a_{r-2},a_{r-3}) - e_{r-4} - \Sig{1}(e_{r-1})
          - \Ch(e_{r-1},e_{r-2},e_{r-3}) - K_r .
\]

\begin{lemma}[Linear entries]
\label{lem:linear}
For every state word $a_k$: $W_k = a_k + f(a_{k-8},\ldots,a_{k-1})$;
$W_{k+8} = -a_k + g(a_{k+1},\ldots,a_{k+8})$; $a_k$ enters
$W_{k+1},\ldots,W_{k+7}$ nonlinearly; and it enters no other $W_r$.
\end{lemma}
\begin{proof}
$W_r$ depends on $a_{r-8},\ldots,a_r$ and nothing else (the nine-word
dependence of~\cite{Basra2026R19}).  In $W_k$ the word $a_k$ occurs once, as
the leading term.  In $W_{k+8}$ it occurs once, inside the $h$-term
$e_{k+4}=a_k+a_{k+4}-T_2(a_{k+3},a_{k+2},a_{k+1})$, which is subtracted.  For
$1\le m\le7$ the dependence is nonlinear: $a_k$ occurs inside a $T_2$, a
$\Sig{1}$ or a $\Ch$ argument, and for $m=4$ additionally as $-a_k$ through
the $h$-term $e_k$, but in no case is the net entry linear.
\end{proof}

The lemma is checked numerically by \texttt{edge\_weights.py}: for $k\le4$
at $R{=}20$ and four random increments each, $W_k$ moves by exactly
$+\delta$, $W_{k+8}$ by exactly $-\delta$, the words between by neither, and
$W_k+W_{k+8}$ is invariant under $a_k$.

\begin{corollary}[Every constraint has a global absorber]
\label{cor:absorb}
$\mathrm{C}_j = W_{16+j}-\sig{1}(W_{14+j})-W_{9+j}-\sig{0}(W_{1+j})-W_j$
contains $a_{j+1}$ exactly twice: as $+a_{j+1}$ inside the $\sig{0}$ argument
$W_{1+j}$ and as $-a_{j+1}$ inside $W_{9+j}$.  With $u=W_{1+j}$ it reads
$\sig{0}(u)-u = (\text{terms free of } a_{j+1})$, so the single global table
absorbs $a_{j+1}$ from $\mathrm{C}_j$ for every $j=0,\ldots,R-17$.
\end{corollary}

For $j=0,1,2$ this is the opposite-sign pattern that drives the $R{=}19$
attack.  For $j=3$ it is C3$\to a_4$.  The inventory of absorptions with
the tables we know is:

\begin{center}
\begin{tabular}{@{}ll@{}}
\toprule
Global (built once, valid for every target and context) & absorbs \\
\midrule
$\sig{0}(u)-u$ & C$_j\to a_{j+1}$ for every $j$: \\
 & C0$\to a_1$, C1$\to a_2$, C2$\to a_3$, C3$\to a_4$ \\
$\Gamma^{-1}$, where $\Gamma(a_0)=a_0+g(a_0)$ (\S\ref{sec:a0absorber}) & C0$\to a_0$ \\
$W_0 = a_0 + \mathrm{const}$ (linear) & --- \\
\midrule
Per-context ($2^{32}$ to build, reusable within one context) & \\
\midrule
$Q_0^{-1},Q_1^{-1},Q_2^{-1},H^{-1}$ & C0--C3 $\to a_{11}$ \\
\midrule
No table (context-dependent inversion, $2^{32}$ every time) & \\
\midrule
C3 absorbing $a_1$, $a_2$ or $a_3$; C0 absorbing $a_4$ & --- \\
\bottomrule
\end{tabular}
\end{center}

\subsection{Two Frames, Both Cyclic}
\label{sec:frames}

Reaching $2^{32}$ per context requires sweeping one unknown and absorbing
four.  Given the tables we know, and up to the $a_0/a_1$ mirror of
\S\ref{sec:a0absorber}, there are exactly two ways to do it at $R{=}20$.
\emph{Frame A} (used in \S\ref{sec:noniter}) keeps $a_4,\ldots,a_{10}$ as
context, sweeps $a_0$, and absorbs $a_1,a_2,a_3$ globally and $a_{11}$
per-context:
\[
  a_1 \leftarrow \mathrm{C0}\ (a_2,a_3,a_{11}), \quad
  a_2 \leftarrow \mathrm{C1}\ (a_1,a_3,a_{11}), \quad
  a_3 \leftarrow \mathrm{C2}\ (a_1,a_2,a_{11}), \quad
  a_{11} \leftarrow \mathrm{C3}\ (a_1,a_2,a_3).
\]
\emph{Frame B} keeps $a_5,\ldots,a_{11}$ as context, sweeps $a_0$, and absorbs
$a_1,\ldots,a_4$ with the one global table:
\[
  a_1 \leftarrow \mathrm{C0}\ (a_2,a_3,a_4), \quad
  a_2 \leftarrow \mathrm{C1}\ (a_1,a_3,a_4), \quad
  a_3 \leftarrow \mathrm{C2}\ (a_1,a_2,a_4), \quad
  a_4 \leftarrow \mathrm{C3}\ (a_1,a_2,a_3).
\]
Sweeping $a_1$ instead of $a_0$ gives the mirror frame of
\S\ref{sec:a0absorber} (sweep $W_1$, absorb $a_0$ through $\Gamma^{-1}$)
and changes nothing below; sweeping $a_2$ or $a_3$ leaves $a_0$ and $a_1$
both to be absorbed by C0, which cannot absorb two words.  Both
assignments are cyclic.  So is the $R{=}19$ assignment
($a_1\leftarrow\mathrm{C0}(a_2,a_3)$, $a_2\leftarrow\mathrm{C1}(a_1,a_3)$,
$a_3\leftarrow\mathrm{C2}(a_1,a_2)$), and it succeeds; cyclicity is not the
barrier.

\subsection{Edge Weights: the Heavy Cycle Appears at \texorpdfstring{$R=20$}{R=20}}
\label{sec:heavy}

What distinguishes the three assignments is how strongly each absorbed word
feeds back into the constraints that absorb the others.  We measure it
directly: for random targets, contexts and unknowns, flip one random bit of
$a_k$ and record the Hamming weight of the change in $\mathrm{C}_j$, averaged
over $200$ trials per cell (\texttt{edge\_weights.py}, seed $20260902$; raw
output in \texttt{data\_edge\_weights.txt}).  Full avalanche is about $16$;
a \emph{mild} edge, one that passes only through $\Maj$ and $\Ch$ of $a_k$
plus a constant, moves one or two bits.  We call a feedback edge
\emph{heavy} when it moves more than $4$ bits; self-edges sit near $8$
because the absorbed word moves its own $\sig{0}(u)-u$ term.

\begin{center}
\small
\begin{tabular}{@{}l rrr@{\hspace{1.6em}} rrrr@{\hspace{1.6em}} rrrrr@{}}
\toprule
 & \multicolumn{3}{c}{$R{=}19$} & \multicolumn{4}{c}{$R{=}20$, frame B} & \multicolumn{5}{c}{$R{=}21$} \\
flip & C0 & C1 & C2 & C0 & C1 & C2 & C3 & C0 & C1 & C2 & C3 & C4 \\
\midrule
$a_0$ (swept) & 15.4 & 15.9 & 15.9 & 15.6 & 15.8 & 15.8 & 15.0 & 15.1 & 16.1 & 15.7 & 15.2 & 15.5 \\
$a_1$ & \underline{8.0} & 15.3 & 16.2 & \underline{7.6} & 15.2 & 16.1 & 15.7 & \underline{8.0} & 15.1 & 16.2 & 15.9 & 15.3 \\
$a_2$ & 1.5 & \underline{7.9} & 15.1 & 1.3 & \underline{7.9} & 14.9 & 16.3 & 1.5 & \underline{7.6} & 14.9 & 16.1 & 16.0 \\
$a_3$ & 1.9 & 1.4 & \underline{8.3} & 1.9 & 1.6 & \underline{8.3} & 15.2 & 1.7 & 1.4 & \underline{7.8} & 15.4 & 16.1 \\
$a_4$ & & & & \textbf{12.2} & 1.7 & 1.4 & \underline{7.8} & \textbf{12.4} & 1.8 & 1.5 & \underline{7.8} & 15.2 \\
$a_5$ & & & & & & & & \textbf{8.3} & \textbf{12.4} & 1.6 & 1.3 & \underline{8.0} \\
\bottomrule
\end{tabular}
\end{center}

Underlined cells are self-edges (the absorbed word moving its own
$\sig{0}(u)-u$ term); entries above the diagonal are forward edges; entries
below it are feedback edges, with the heavy ones in bold.  Frame A at
$R{=}20$ (rows $a_1,a_2,a_3,a_{11}$, same script) has feedback weights
$a_{11}\to$C0 $15.7$, $a_{11}\to$C1 $13.3$, $a_{11}\to$C2 $15.0$.

\paragraph{Reading.}
At $R{=}19$ every feedback edge is mild.  The full-avalanche (``heavy'')
dependency graph is the chain $a_1\to a_2\to a_3$, which is acyclic, and the
attack exploits exactly this: it computes $a_1$ outside the loop, iterates
$a_3\to a_2\to a_3'$, and thereby cuts its cycle at the mild edge
$a_3\to$C1.  That is why the iterated map is nothing like a random map (its
convergence rate of $5.34\times10^{-5}$ sits $2^{15}$ above the random-map
baseline, \S\ref{sec:noniter}).  At $R{=}20$ exactly one feedback edge is
heavy, $a_4\to$C0 at $12.2$ bits, and the heavy graph acquires the cycle
$a_1\to a_2\to a_3\to a_4\to a_1$.  A heavy cycle has no mild cut.  Cutting
it at any single $32$-bit word composes full-avalanche maps with table
lookups, which behaves, as far as we can measure, like a random map on
$2^{32}$ points; its fixed points then cost $2^{32}$ each to find, which is
the classic sweep.  Frame A is worse, with three heavy feedback edges
through $a_{11}$.  This is the content of the ``dimensional'' explanation
in \S\ref{sec:noniter}: the state cannot be held at $32$ bits because no
mild cut exists.

\paragraph{Why it is the round function and not the frame.}
The weight of an entry is set by its offset.  $a_k$ enters $W_{k+m}$ for
$m=6,7$ only through $\Maj$ and $\Ch$ of $a_k$ plus constants (mild); for
$m=5$ through the $h$-term $e_{k+1}$, which contains $\Sig{0}(a_k)$ subtracted
directly, and through $\Sig{1}(e_{k+4})=\Sig{1}(a_k+c)$ (heavy; measured
$12.2$--$12.4$); for $m=4$ through $\Ch(e_{k+3},e_{k+2},e_{k+1})$, which
selects about half the bits of $\Sig{0}(a_k)$ (measured $8.3$); and for
$m\le3$ through $\Sig{0}$ and $\Sig{1}$ of $a_k$ itself.  The last word of
the window, $a_n$ with $n=R-16$, feeds C0 through $W_9$ at offset $9-n$, which
is mild if and only if $n\le3$, i.e.\ $R\le19$.  No context choice removes
$\Sig{0}(a_4)$ from $W_9$: it enters additively with coefficient $+1$
through the subtracted $h$-term $e_5$, and $a_4$ reaches C0 through no
other word, since $W_0$, $W_1$, $W_{14}$ and $W_{16}$ all lie outside
offsets $0$--$8$ from $a_4$.  At
$R{=}21$ the window has three heavy feedback edges (table), and the count
grows from there.

\paragraph{What would cut it.}
A trick that cancels the barrier must cut one heavy edge, and the cheapest
candidate is the new one: in C0, invert $\Sig{0}(a_4)-\Sig{1}(a_4+c_8)$ plus
its $\Maj/\Ch$ terms \emph{jointly} with the $a_1$ absorber $\sig{0}(W_1)-W_1$
at well below $2^{32}$ per instance.  All of these are rotation--XOR maps, so
over $\mathrm{GF}(2)$ they would combine linearly; the modular carries between
them are what destroy the combination, and \S\ref{sec:leftnull} measures that
destruction for $\sig{0}$ alone (an exact $\mathrm{GF}(2)$ identity degraded
to a $2^{-15.8}$ bias).  Other routes one might try---a normal form for
$\sig{0}(W_4)+W_3$ in one message word, freeing the sweep variable through
the $a_0$ absorber, a joint two-variable global table---are all instances
of cutting a heavy edge, and none has been found.
The two attempts recorded below are of this kind and both charge full price.

\subsection{Attempts to Cut the Cycle}
\label{sec:cut}

\paragraph{The one structural handle, and why it does not pay.}
C3 has the form $\sig{0}(W_4)+W_3 = H(a_{11})$, coupling \emph{two} message
words---this is precisely why it cannot absorb $a_1,a_2,a_3$ through the
single-word table.  Imposing any global relation $W_3=\phi(W_4)$ collapses the
left side to a function of one word, making C3 a global absorber producing
$(W_3,W_4)$ from $a_{11}$ alone and reversing the cycle direction.  Measured
exhaustively in the reduced-width model at $w=4$ over 6 contexts (98 true
solutions; expectation $6.12$ if $\phi$ costs a full $w$ bits):

\begin{center}
\begin{tabular}{@{}lrrr@{}}
\toprule
Relation & surviving & per context & vs.\ expected \\
\midrule
$W_3=\sig{0}(W_4)$ & 3 & 0.50 & $0.49\times$ \\
$W_3=-W_4$         & 4 & 0.67 & $0.65\times$ \\
$W_3=W_4$          & 0 & 0.00 & $0.00\times$ \\
$W_3=\overline{W_4}$ & 7 & 1.17 & $1.14\times$ \\
\bottomrule
\end{tabular}
\end{center}

Every variant sits at or near the ``costs a full $w$ bits'' expectation.  The
counts are small ($3$--$7$ events, so roughly $\pm2$) and all are consistent
with $1.0\times$; none is consistent with the large excess needed to make the
handle pay.  Imposing a $w$-bit condition removes a $w$-bit worth of solutions,
and the global absorber it buys does not compensate.

\paragraph{Elimination orders.}
An elimination order cannot change an edge weight, so the heavy-cycle
argument predicts that no order recovers the solution set.  The exhaustive
search over all $2{,}880$ elimination orders in the reduced-width model
($w{=}4$, two contexts, $41$ true solutions) agrees: $2{,}560$ of
$2{,}880$ orders find at least one solution, the median finds $2$ and the
best finds $11$, and the hand-chosen $R{=}19$-style order (sweep $a_0$;
$a_1\leftarrow$C0, $a_2\leftarrow$C1, $a_3\leftarrow$C2,
$a_{11}\leftarrow$C3) ranks $175$th with $5$.  No order recovers more than
$27\%$ of the solution set.  The search treats every absorption as
available, ignoring the table inventory, so it is if anything generous to
the alternatives.

\section{Status}

$R=20$ is not solved.  The barrier is $2^{64}$ per context-equivalent for
this attack family, $2^{32}$ times the $R{=}19$ per-context cost, or about
$2^{71.4}$ expected total work at the projected $R{=}19$ rate, and its
sharpest form is the heavy-cycle argument of \S\ref{sec:heavy}: the
full-avalanche dependency graph among the absorbed words is acyclic for
$R\le19$ and cyclic from $R{=}20$ in every admissible frame, because
$\Sig{0}(a_k)$ is subtracted directly into $W_{k+5}$.  Context screening
(\cite[\S5.7]{Basra2026R19}) raises $W_9$ lo-pass yield by a measured
$3.41\times$ ($95\%$ CI $[2.13,5.82]$) with conversion to preimages
unmeasured; taken at face value that lowers the expected total work by the
same factor, to about $2^{69.6}$, an improvement and not a solution.

We emphasise the limits of the negative results.  Each rules out a
parameterisation, not a technique: the $\phi$-family charges a full word for
the global absorber it buys; $a_0$ is globally absorbable but competes with
$a_1$ for C0; no elimination order recovers more than a quarter of the
solution set.  The heavy-cycle
argument is a measurement of this attack family, not a lower bound.  What
would change the picture is a way to invert $\Sig{0}(a_4)-\Sig{1}(a_4+c_8)$
jointly with $\sig{0}(W_1)-W_1$ inside C0 at well below $2^{32}$, and nothing
measured here excludes one; it does say that any such trick is a statement
about carries between rotation--XOR maps, which is where \S\ref{sec:leftnull}
found only a $2^{-15.8}$ bias.

\begin{thebibliography}{10}

\bibitem{Basra2026R19}
K.~S.~Basra.
\newblock A practical preimage attack on the 19-round {SHA-256} compression
  function via a global $\sigma_0$-difference table.
\newblock Main paper, submitted concurrently, 2026.  Source and code at
  \url{https://github.com/kamb-code/sha256-r19-preimage}.

\bibitem{AokiGuo2009}
K.~Aoki, J.~Guo, K.~Matusiewicz, Y.~Sasaki, and L.~Wang.
\newblock Preimages for step-reduced {SHA-2}.
\newblock In \emph{Advances in Cryptology -- ASIACRYPT 2009}, LNCS 5912,
  pp.~578--597. Springer, 2009.

\bibitem{KRS2012}
D.~Khovratovich, C.~Rechberger, and A.~Savelieva.
\newblock Bicliques for preimages: attacks on {Skein-512} and the {SHA-2}
  family.
\newblock In \emph{Fast Software Encryption (FSE 2012)}, LNCS 7549,
  pp.~244--263. Springer, 2012.

\bibitem{Zaikin2026}
O.~Zaikin.
\newblock Preimage attacks on round-reduced {MD5}, {SHA-1}, and {SHA-256}
  using parameterized {SAT} solver.
\newblock \emph{Constraints}, published online 23 January 2026.
  DOI 10.1007/s10601-025-09383-0.

\end{thebibliography}

\end{document}
